\section*{AI Assistance}
\label{sec:ai-assistance}

We used LLM-based assistance, including OpenAI Codex, to draft, restructure, and edit the manuscript; prepare a Chinese reading version; search for related literature and bibliographic information; and assist with analysis and plotting code. Generative image tools were used to create and revise explanatory framework figures.

AI assistance also contributed to the information-theoretic interpretation of the architecture, the formulation of mathematical claims and their supporting arguments, and the discussion of experimental methodology and results. This included the treatment of progressive entity alignment, search--evidence separation, and query-dependent reading budgets. These arguments are conditional on the assumptions stated in the paper; AI assistance does not constitute independent proof verification or empirical validation.

No synthetic evaluation dataset was created as part of this paper-preparation workflow. Reported benchmark scores are drawn from recorded experimental runs rather than generated performance estimates. The authors remain responsible for the final text, citations, mathematical claims, code, figures, and empirical conclusions, including verification of AI-assisted material.
